% Technical report -- skeleton (2026-09-14, night build; Fable plan item D).
% Every number here is copied from a result file or from docs/project_notes.tex and is
% keyed to its tag. Cells marked pending are filled from benchmark/paired_stats.json and
% benchmark/protocol_results_<tag>.json once those jobs land. Nothing here is invented;
% where a number does not exist yet the cell says so.
\documentclass[11pt,a4paper]{article}
\usepackage[utf8]{inputenc}
\usepackage[T1]{fontenc}
\usepackage{lmodern}
\usepackage{microtype}
\usepackage{amsmath}
\usepackage[margin=1in]{geometry}
\usepackage{booktabs}
\usepackage{graphicx}
\usepackage{xcolor}
\usepackage[hidelinks]{hyperref}
\newcommand{\todo}[1]{\textcolor{red}{[#1]}}
\newcommand{\tagref}[1]{\texttt{#1}}

\title{Visual Augmentation of Audio Semantics for Accessibility\\
\large Technical report}
\author{Adam Gavriely\\Bar-Ilan University, Department of Computer Science}
\date{September 2026}

\begin{document}
\maketitle

\begin{abstract}
\todo{Write last. One paragraph: the task (show a deaf viewer the sounds whose source is not on
screen), the pipeline (seven stages of off-the-shelf models, no training), the benchmark (274
clips, four categories), the evaluation (VLM--LLM protocol with an independent reference), the
headline (\tagref{v3\_indep} table), and the two findings that survive every reference: the gate is
bounded by what a VLM can answer about visibility, and the automatic reference must be built by
models that never saw the system's output.}
\end{abstract}

% ======================================================================
\section{Introduction}
% Motivation, problem statement, the accessibility gap (subtitles/SDH and sign language
% carry speech; ambient sound is described inconsistently or not at all).
% Source: project_notes.tex sec:gap.
\todo{From notes \S{}The Accessibility Gap.}

\subsection{Problem statement}
Given a video, produce the same video with a side panel that shows, for the duration of each
ambient sound whose source is not visible on screen, a picture of that sound's event -- and shows
nothing otherwise. Speech and music are excluded. No model is trained; the system is an inference
pipeline over open-source foundation models, and human involvement is limited to the benchmark
labels used for evaluation.

\subsection{Research questions}
% Verbatim from the proposal (Aug 6, 2026); each is answered in sec:discussion.
\begin{enumerate}
  \item Which types of audio information contribute most to scene understanding when presented visually?
  \item What level of semantic granularity should be used to represent ambient sounds?
  \item Can automatically generated visual augmentations improve accessibility beyond subtitles?
  \item How should semantic audio augmentation systems be evaluated?
  \item How does the proposed pipeline compare with representative audio-to-visual approaches?
\end{enumerate}

\subsection{Contributions}
\begin{itemize}
  \item Formulation of semantic audio augmentation as an accessibility task, with a precise
        definition of when a picture is due (sound heard, source not visible, not speech or music).
  \item A seven-stage inference pipeline of off-the-shelf models (\S\ref{sec:method}).
  \item A benchmark of 274 human-labelled clips in four scenarios (\S\ref{sec:benchmark}).
  \item An automatic evaluation protocol, and the finding that its reference must be independent
        of the system under test (\S\ref{sec:eval}).
  \item An experimental comparison with two baselines, with paired statistics and an oracle bound
        (\S\ref{sec:results}).
\end{itemize}

% ======================================================================
\section{Related work}
% Condense literature_review.tex: AED/tagging (PANNs, BEATs), audio captioning, audio LLMs
% (Qwen2-Audio), VLMs (Qwen2.5-VL, Idefics3), audio-to-image (Sound2Scene etc.) as the
% baseline family, SDH/sign-language handling of non-speech sound, DHH visualisation studies.
\todo{Condense \texttt{docs/literature\_review.tex}; one paragraph per area, positioned against
this project.}

% ======================================================================
\section{Method}\label{sec:method}

\subsection{Pipeline}
% Table: stage, model, role. Source: README pipeline table; config.py.
\begin{center}\begin{tabular}{lll}
\toprule
stage & model & role \\
\midrule
1 audio extraction & FFmpeg & 16 kHz mono wav \\
2 video understanding & OWLv2 (cheap pass) & candidate visible sources, deferred to stage 5 \\
3 speech recognition & Whisper (configurable) & transcript, used only as gate context \\
4 audio event detection & BEATs iter3+ AS2M (527 AudioSet classes) & events with onsets \\
5 cross-modal reasoning & Qwen2.5-VL-7B & place, visibility, kind, depiction, dedup \\
6 visual augmentation & FLUX.1-schnell, 768 px & one picture per sound, side panel \\
7 evaluation & Qwen2.5-VL / Mistral-7B judge + independent reference & \S\ref{sec:eval} \\
\bottomrule
\end{tabular}\end{center}

\subsection{Stage 4: detection and onset timing}
% BEATs replaced PANNs (notes sec:beats: 11/14 detector complaints resolved). 2 s window,
% hop 0.25 s, reflect lead-in pad, stamp offset 0.5 s. Threshold 0.35, set once (plan_robustness).
% Onset: hysteresis (0.5 of threshold) + prefix-silencing occlusion (notes sec:onset).
% CAM on BEATs tokens does not work (deltas exactly 0 / -1.5): recorded as a negative result.
\todo{From notes \S{}The detector was the bottleneck and \S{}When does a sound start.}

\subsection{Stage 5: is the source on screen?}
% Majority of three signals (open naming, world-knowledge check, open description), every
% two-way question asked in both orderings because the 7B model has a position bias; judged
% per 5 s stretch; one continuous picture per continuous sound (notes sec:vlmvisibility,
% sec:actionvisible). Speech as gate context, never as picture material (sec:speechgate).
\todo{From notes \S{}Visibility was a list of thirty things, \S{}Visible means the action,
\S{}Speech as gate context.}

\subsection{Stage 5: one picture per source}
% Ontology dedup (AudioSet ancestor rule; family vs member by interval subtraction; similarity
% merges need time overlap) -- three attempts recorded (notes sec:dedup). Disambiguation among the
% detector's own candidates from frames (sec:disamb).
\todo{From notes \S{}One picture per source.}

\subsection{Stage 6: what the picture shows}
% Event-framed depiction (2-5 plain words, no place, no example sentences), forced-choice
% validation asked twice; FLUX over PixArt; panel with dwell, bursts, three rows.
\todo{From notes \S{}A place not a scene sentence, \S{}Validating a depiction, \S{}The panel was a collage.}

\subsection{Design process}
% Knobs set once on the dev split; failure catalogue (18 entries) with the >=3-cases rule;
% demos illustrate only. docs/plan_robustness.md.
\todo{Summarise \texttt{docs/plan\_robustness.md}.}

% ======================================================================
\section{Benchmark}\label{sec:benchmark}
% 274 clips (proposal named 300; usable yield 15-20%), four categories: unseen_ambient, mixed,
% seen_ambient, no_ambient. One annotator; blind re-adjudication kappa = 0.60 (notes
% sec:readjudication). First 209 tags predate any model filtering. Main runs use 100 clips,
% 25 per category.
\todo{From notes \S{}Sourcing yield, \S{}Annotation guidelines, \S{}Label stability.}

\subsection{External check: DCASE 2025 Task 3 on/off-screen labels}
% benchmark/eval_dcase_visibility.json, 258 events, gold on/off-screen from the dataset.
\begin{center}\begin{tabular}{lc}
\toprule
visibility check vs.\ DCASE 2025 gold (258 events) & \\
\midrule
agreement & 66.7\% \\
off-screen recall (sound correctly not silenced) & 97\% \\
on-screen recall (visible source correctly silenced) & 35\% \\
\bottomrule
\end{tabular}\end{center}
The check is timid in one direction: it almost never silences a sound whose source is off screen,
and silences only a third of those whose source is on screen. \todo{DCASE geometric caveat:
360-degree capture; on-screen is defined by the field of view of the equirectangular frame.}

% ======================================================================
\section{Evaluation protocol}\label{sec:eval}
% Proposal 6.1: VLM describes the augmented scene -> LLM writes a reference -> independent judge
% scores 0-4. Abstention scoring (notes sec:abstain); reference must be able to say "nothing is
% missing" (sec:reference).
\subsection{Three references}
\begin{description}
  \item[model-derived (\tagref{v2}, \tagref{v3})] as specified in the proposal: written by the
        system's own VLM from the detector's own events. Found circular on 2026-09-14 (review): it
        rewards whichever system repeats the detector.
  \item[human-grounded (\tagref{v2\_grounded}, \tagref{v3\_grounded})] the same reference corrected
        by the annotator's label; only computable on labelled clips; raises the abstaining system
        more than the baselines by arithmetic (notes \S{}Main Result).
  \item[independent (\tagref{v2\_indep}, \tagref{v3\_indep})] built by models that never see the
        system's output: Qwen2-Audio listens to the Demucs no-vocals stem in 10 s chunks; each claim
        is verified by CLAP against 64 decoys (mean $+2\sigma$); Idefics3-8B looks at the frames for
        what is on screen; Llama-3.1-8B writes the reference; a MiniLM sentinel ($\tau = 0.50$,
        uncalibrated) rejects sentences unsupported by the verified list. This is the corrected
        primary.
\end{description}

\subsection{Baselines}
Blind audio-to-image (depict everything heard, no gate) and audio captioning (text instead of a
picture); video reasoning off for both.

\subsection{Statistics}
Per-clip paired win/tie/loss, bootstrap 95\% CI on the mean difference (5000 resamples), and an
oracle-gate bound: the human tag as the gate, the best system's output where a picture is due
(\texttt{scripts/paired\_stats.py}).

% ======================================================================
\section{Results}\label{sec:results}

\subsection{Main comparison}
% Fill from benchmark/paired_stats.json. v2 = 2026-09-12 configuration (PANNs, PixArt, 30-concept
% visibility list). v3 = shipping configuration (BEATs, FLUX, per-sound visibility, ontology dedup).
% Scores 0-4, 100 clips, 25 per category, Mistral-7B judge.
\begin{center}\begin{tabular}{llccc}
\toprule
run & reference & proposed & blind a2i & audio caption \\
\midrule
\tagref{v2} & model-derived   & 2.27 & \textbf{2.97} & 2.68 \\
\tagref{v2\_grounded} & human-grounded & 2.90 & \textbf{3.14} & 2.79 \\
\tagref{v2\_indep} & independent & \todo{pending job 28942612} & & \\
\midrule
\tagref{v3} & model-derived   & \todo{pending job 28942456} & & \\
\tagref{v3\_grounded} & human-grounded & \todo{pending} & & \\
\tagref{v3\_indep} & independent & \todo{after v3 describe: reference phase, then judge} & & \\
\bottomrule
\end{tabular}\end{center}

\subsection{Paired statistics (\tagref{v2\_grounded})}
% scripts/paired_stats.py v2_grounded, 2026-09-14.
proposed vs.\ blind: win 22 / tie 54 / loss 24, mean difference $-0.24$, 95\% CI $[-0.47, -0.02]$.
Oracle gate + best output: 3.19. \todo{Add the same line for \tagref{v2\_indep} and \tagref{v3\_indep}.}

\subsection{Per-category}
\todo{Per human tag table from \texttt{paired\_stats.py} output, per reference.}

\subsection{Onset timing}
% scripts/timing_compare.py; 7 hand-labelled onsets on 3 clips (sizzling 0, alarm ~9, siren ~9,
% horse 17.5, dog 0.5, helicopter ~8.5, door ~1.5). n = 7: indicative only.
Mean onset error against seven hand-labelled onsets on three clips: sliding-window stamp $+0.9$ s,
hysteresis + occlusion $+0.3$ s ($n = 7$; a DCASE-scale onset evaluation is listed as future work,
plan item C).

\subsection{Ablations recorded during development}
% From the notes: retrieval vs SDXL (+0.17/+0.18, earlier configuration); judge reliability
% (kappa 0.753 second judge); PANNs -> BEATs (11/14 complaints resolved); CAM negative result;
% place veto 2 right / 2 wrong on 4 uses -> off.
\todo{One short paragraph each, all numbers from the notes.}

% ======================================================================
\section{Discussion}\label{sec:discussion}
% Answer the five research questions with the numbers above. Be plain where the gate does not win.
\todo{RQ1--RQ5.}

% ======================================================================
\section{Limitations}
% Mirror LIMITATIONS.md (repository root); keep the two in step.
\todo{From \texttt{LIMITATIONS.md}.}

% ======================================================================
\section{Conclusion and future work}
\todo{Open-vocabulary audio (CLAP zero-shot / audio LLM) as a second opinion; source-separation front
end; human study; second annotator; $\tau$ calibration of the sentinel on DCASE gold.}

\bibliographystyle{plain}
\bibliography{../project_notes}
\end{document}
